The inflation we observed follows from the pooled variance. With 40 observations at standard deviation 1 and 10 at standard deviation 3, the pooled estimate sits much closer to 1 than to 3, so the standard error of the difference is too small and the $t$ statistic too large. Welch's test keeps the two variances apart and avoids this.

The simulation covers one variance ratio, one pair of sample sizes and normal data. Errors would shrink as the groups become more similar in size, and the direction reverses when the larger group has the larger variance, where Student's test becomes conservative and loses power~\citep{ruxton2006unequal}. Heavy-tailed data would add a separate problem that neither test addresses. A useful follow-up would vary the size ratio from 1 to 8 and the standard deviation ratio from 1 to 4 and report where the error rate of Student's test first exceeds 7.5\%, the upper bound often used to judge robustness.
